\documentclass[10pt,a4paper]{amsart}
\usepackage{amsmath,amssymb,amsthm,graphicx}
\usepackage{verbatim}
\newenvironment{reference}{\comment}{\endcomment}
%\newenvironment{reference}{}{}
\newenvironment{slogan}{\comment}{\endcomment}
\newenvironment{history}{\comment}{\endcomment}
% For commutative diagrams we use Xy-pic
\usepackage[all]{xy}

% We use 2cell for 2-commutative diagrams.
\xyoption{2cell}
% Theorem environments.
%
\theoremstyle{plain}
\newtheorem{theorem}[subsection]{Theorem}
\newtheorem{proposition}[subsection]{Proposition}
\newtheorem{lemma}[subsection]{Lemma}

\theoremstyle{definition}
\newtheorem{definition}[subsection]{Definition}
\newtheorem{example}[subsection]{Example}
\newtheorem{exercise}[subsection]{Exercise}
\newtheorem{situation}[subsection]{Situation}

\theoremstyle{remark}
\newtheorem{remark}[subsection]{Remark}
\newtheorem{remarks}[subsection]{Remarks}

\numberwithin{equation}{subsection}

% Macros
%
\def\lim{\mathop{\mathrm{lim}}\nolimits}
\def\colim{\mathop{\mathrm{colim}}\nolimits}
\def\Spec{\mathop{\mathrm{Spec}}}
\def\Hom{\mathop{\mathrm{Hom}}\nolimits}
\def\Ext{\mathop{\mathrm{Ext}}\nolimits}
\def\SheafHom{\mathop{\mathcal{H}\!\mathit{om}}\nolimits}
\def\SheafExt{\mathop{\mathcal{E}\!\mathit{xt}}\nolimits}
\def\Sch{\mathit{Sch}}
\def\Mor{\mathop{\mathrm{Mor}}\nolimits}
\def\Ob{\mathop{\mathrm{Ob}}\nolimits}
\def\Sh{\mathop{\mathit{Sh}}\nolimits}
\def\NL{\mathop{N\!L}\nolimits}
\def\CH{\mathop{\mathrm{CH}}\nolimits}
\def\proetale{{pro\text{-}\acute{e}tale}}
\def\etale{{\acute{e}tale}}
\def\QCoh{\mathit{QCoh}}
\def\Ker{\mathop{\mathrm{Ker}}}
\def\Im{\mathop{\mathrm{Im}}}
\def\Coker{\mathop{\mathrm{Coker}}}
\def\Coim{\mathop{\mathrm{Coim}}}

% Boxtimes
%
\DeclareMathSymbol{\boxtimes}{\mathbin}{AMSa}{"02}

%
% Macros for moduli stacks/spaces
%
\def\QCohstack{\mathcal{QC}\!\mathit{oh}}
\def\Cohstack{\mathcal{C}\!\mathit{oh}}
\def\Spacesstack{\mathcal{S}\!\mathit{paces}}
\def\Quotfunctor{\mathrm{Quot}}
\def\Hilbfunctor{\mathrm{Hilb}}
\def\Curvesstack{\mathcal{C}\!\mathit{urves}}
\def\Polarizedstack{\mathcal{P}\!\mathit{olarized}}
\def\Complexesstack{\mathcal{C}\!\mathit{omplexes}}
% \Pic is the operator that assigns to X its picard group, usage \Pic(X)
% \Picardstack_{X/B} denotes the Picard stack of X over B
% \Picardfunctor_{X/B} denotes the Picard functor of X over B
\def\Pic{\mathop{\mathrm{Pic}}\nolimits}
\def\Picardstack{\mathcal{P}\!\mathit{ic}}
\def\Picardfunctor{\mathrm{Pic}}
\def\Deformationcategory{\mathcal{D}\!\mathit{ef}}
\usepackage{cleveref}
\crefname{originalreference}{original reference}{original references}
\Crefname{originalreference}{Original reference}{Original references}
\makeatletter
\newlabel{categories-definition-dual}{{categories-definition-dual}{source}{}{sourcepointer-1}{}}
\newlabel{categories-definition-dual@cref}{{[originalreference][0][]categories-definition-dual}{[1][0][]source}{}{}{}}
\makeatother
\crefname{originalreference}{original reference}{original references}
\Crefname{originalreference}{Original reference}{Original references}
\makeatletter
\newlabel{categories-lemma-left-dual}{{categories-lemma-left-dual}{source}{}{sourcepointer-2}{}}
\newlabel{categories-lemma-left-dual@cref}{{[originalreference][0][]categories-lemma-left-dual}{[1][0][]source}{}{}{}}
\makeatother
\crefname{originalreference}{original reference}{original references}
\Crefname{originalreference}{Original reference}{Original references}
\makeatletter
\newlabel{categories-remark-left-dual-adjoint}{{categories-remark-left-dual-adjoint}{source}{}{sourcepointer-3}{}}
\newlabel{categories-remark-left-dual-adjoint@cref}{{[originalreference][0][]categories-remark-left-dual-adjoint}{[1][0][]source}{}{}{}}
\makeatother
\crefname{originalsection}{original section}{original sections}
\Crefname{originalsection}{Original section}{Original sections}
\makeatletter
\newlabel{example-dual}{{example-dual}{source}{}{sourcepointer-4}{}}
\newlabel{example-dual@cref}{{[originalsection][0][]example-dual}{[1][0][]source}{}{}{}}
\makeatother
\title{A module admitting a left dual is locally a direct summand of a finite free module}
\author{The Stacks Project Authors}
\date{stacks-project:modules.tex:local-snapshot-20261012; GNU FDL 1.2 or later}
\begin{document}
\maketitle
\noindent\textit{Editorial note.} All delivered mathematical text is the author's own, copied from the pinned Stacks Project chapter source sources/stacks-modules/pinned-original-source.tex; no statement or proof step was rewritten, reordered or removed.\par\medskip
\noindent\textit{Editorial note.} Preamble provenance: the Stacks Project keeps one single shared preamble for every chapter; this card takes it directly from the staged sources/stacks-modules/preamble.tex as three verbatim spans (lines 8-12, 14-18 and 154-217), so no derived file was written and nothing was added or redefined. The wrapper class and the base packages are supplied by the assembler.\par\medskip
\noindent\textit{Editorial note.} Bibliography provenance: neither the counted statement nor its proof contains a citation (this lane rescans the whole span before writing the card), so the card carries no bibliography and no entry was reconstructed.\par\medskip
\noindent\textit{Editorial note.} Selection provenance: this card was opened under the file-based control contract (CONTROL.md), with the mechanical pre-card gate run on its statement and proof interval before the card was written.\par\medskip
\noindent\textit{Allowed standard prerequisites (editorial).}\begin{itemize}
\item The notion of a left dual of an object of a monoidal category, witnessed by a unit and a counit whose composites with the coherence maps are identities. Invoked in the statement (source lines 3089-3091) and in the counted proof (source lines 3105-3116) as a cross-chapter pointer. Supplies the two identities that the proof unfolds at the start.
\item The description of the unit as a sum of finitely many products of local sections near a point, and the adjunction between morphisms of modules. Used inside the counted proof at source lines 3117-3128. Produces the finite free module through which the identity factors locally.
\end{itemize}\medskip
\ifdefined\hypertarget\hypertarget{sourcepointer-1}{}\fi
\noindent\textit{Source reference.} Original reference categories-definition-dual: sources/stacks-modules/pinned-original-source.tex, source lines 3077--3077. Navigation pointer only: the counted text refers to the Stacks item categories-definition-dual, which belongs to another chapter of the project and is not part of this card. The pointer carries the target's own label as its number because this chapter's staged snapshot does not contain the other chapter's tag table; its content is not reproduced here.\par\medskip
\ifdefined\hypertarget\hypertarget{sourcepointer-2}{}\fi
\noindent\textit{Source reference.} Original reference categories-lemma-left-dual: sources/stacks-modules/pinned-original-source.tex, source lines 3145--3145. Navigation pointer only: the counted text refers to the Stacks item categories-lemma-left-dual, which belongs to another chapter of the project and is not part of this card. The pointer carries the target's own label as its number because this chapter's staged snapshot does not contain the other chapter's tag table; its content is not reproduced here.\par\medskip
\ifdefined\hypertarget\hypertarget{sourcepointer-3}{}\fi
\noindent\textit{Source reference.} Original reference categories-remark-left-dual-adjoint: sources/stacks-modules/pinned-original-source.tex, source lines 3148--3148. Navigation pointer only: the counted text refers to the Stacks item categories-remark-left-dual-adjoint, which belongs to another chapter of the project and is not part of this card. The pointer carries the target's own label as its number because this chapter's staged snapshot does not contain the other chapter's tag table; its content is not reproduced here.\par\medskip
\ifdefined\hypertarget\hypertarget{sourcepointer-4}{}\fi
\noindent\textit{Source reference.} Original section example-dual: sources/stacks-modules/pinned-original-source.tex, source lines 3039--3039. Navigation pointer only: the counted text refers to the Stacks item example-dual, which belongs to another chapter of the project and is not part of this card. The pointer carries the target's own label as its number because this chapter's staged snapshot does not contain the other chapter's tag table; its content is not reproduced here.\par\medskip
\begin{lemma}
\label{lemma-left-dual-module}
Let $(X, \mathcal{O}_X)$ be a ringed space. Let $\mathcal{F}$ be an
$\mathcal{O}_X$-module. Let $\mathcal{G}, \eta, \epsilon$
be a left dual of $\mathcal{F}$ in the monoidal category of
$\mathcal{O}_X$-modules, see
Categories, Definition \ref{categories-definition-dual}. Then
\begin{enumerate}
\item $\mathcal{F}$ is locally a direct summand of a finite free
$\mathcal{O}_X$-module,
\item the map
$e : \SheafHom_{\mathcal{O}_X}(\mathcal{F}, \mathcal{O}_X) \to \mathcal{G}$
sending a local section $\lambda$ to $(\lambda \otimes 1)(\eta)$
is an isomorphism,
\item we have $\epsilon(f, g) = e^{-1}(g)(f)$ for local sections
$f$ and $g$ of $\mathcal{F}$ and $\mathcal{G}$.
\end{enumerate}
\end{lemma}

\begin{proof}
The assumptions mean that
$$
\mathcal{F} \xrightarrow{\eta \otimes 1}
\mathcal{F} \otimes_{\mathcal{O}_X} \mathcal{G}
\otimes_{\mathcal{O}_X} \mathcal{F}
\xrightarrow{1 \otimes \epsilon} \mathcal{F}
\quad\text{and}\quad
\mathcal{G} \xrightarrow{1 \otimes \eta}
\mathcal{G} \otimes_{\mathcal{O}_X} \mathcal{F}
\otimes_{\mathcal{O}_X} \mathcal{G}
\xrightarrow{\epsilon \otimes 1} \mathcal{G}
$$
are the identity map. Let $x \in X$. We can find an open neighbourhood
$U$ of $x$, a finite number of sections $f_1, \ldots, f_n$
and $g_1, \ldots, g_n$ of
$\mathcal{F}$ and $\mathcal{G}$ over $U$ such that
$\eta(1) = \sum f_i g_i$. Denote
$$
\mathcal{O}_U^{\oplus n} \to \mathcal{F}|_U
$$
the map sending the $i$th basis vector to $f_i$. Then we
can factor the map $\eta|_U$ over a map
$\tilde \eta : \mathcal{O}_U \to
\mathcal{O}_U^{\oplus n} \otimes_{\mathcal{O}_U} \mathcal{G}|_U$.
We obtain a commutative diagram
$$
\xymatrix{
\mathcal{F}|_U
\ar[rr]_-{\eta \otimes 1} \ar[rrd]_-{\tilde \eta \otimes 1} & &
\mathcal{F}|_U \otimes \mathcal{G}|_U \otimes \mathcal{F}|_U
\ar[r]_-{1 \otimes \epsilon} &
\mathcal{F}|_U \\
& &
\mathcal{O}_U^{\oplus n} \otimes \mathcal{G}|_U \otimes \mathcal{F}|_U
\ar[u] \ar[r]^-{1 \otimes \epsilon} &
\mathcal{O}_U^{\oplus n} \ar[u]
}
$$
This shows that the identity on $\mathcal{F}$ locally on $X$
factors through a finite free module. This proves (1). Part (2) follows from
Categories, Lemma \ref{categories-lemma-left-dual} and its proof.
Part (3) follows from the first equality of the proof.
You can also deduce (2) and (3) from the uniqueness of left duals
(Categories, Remark \ref{categories-remark-left-dual-adjoint})
and the construction of the left dual in
Example \ref{example-dual}.
\end{proof}

\end{document}
